\begin{comment}
{\color{MidnightBlue}
\vspace{5mm}

\textbf{Conjuntos e índices:}
\begin{align*}
I &: \text{Instrumentos posibles} \,(i \in I).\\
O &: \text{Observaciones posibles} \, (o \in O).
\end{align*}

\textbf{Parámetros:}
\begin{align*}
M &: \text{Cantidad máxima de instrumentos a instalar.}\\
\bar{t}_i &: \text{Tiempo máximo de operación del instrumento } i.\\
f_i &: \text{Costo fijo de utilizar el instrumento } i.\\
v_i &: \text{Costo variable por cada hora de operación del instrumento.}\\
B &: \text{Presupuesto de operación.}\\
a_{o,i} &: \text{Parámetro que indica si la observación } o \text{ necesita efectuarse con el instrumento } i.\\
t_{o,i} &: \text{Tiempo de exposición del instrumento } i \text{ necesario para realizar observación } o.\\
u_o &: \text{Utilidad científica de realizar la observación } o.
\end{align*}

\textbf{Variables:}
\begin{align*}
y_{o,i} &\in \{0,1\}: \text{ 1 si el instrumento } i \text{ se utiliza para la observación } o,\text{ 0 en caso contrario.}\\
x_i &\in \{0,1\}: \text{ 1 si el instrumento } i \text{ se instala antes del inicio de la jornada, 0 en caso contrario.}
\end{align*}

\textbf{Función objetivo:}
\begin{equation}
\max_{x_i, y_{o,i}} \quad \sum_{i\in I}\sum_{o\in O} u_o\,y_{o,i}
\end{equation}

\textbf{Restricciones:}

\begin{itemize}[label=\textcolor{MidnightBlue}{\textbullet}]
    \item Máxima cantidad de instrumentos instalados simultáneamente.
    \begin{equation}
        \sum_{i\in I} x_i \leq M
    \end{equation}
    \item Máximo tiempo de operación por instrumento.
    \begin{equation}
        \sum_{o\in O} t_{o,i}\,y_{o,i} \leq \bar{t}_i\,x_i \quad \forall i\in I
    \end{equation}
    \item El instrumento debe haberse instalado para poder ser utilizado.
    \begin{equation}
        y_{o,i} \leq x_i \quad \forall i\in I,\ \forall o\in O
    \end{equation}
    \item Cada observación se puede realizar a lo más una vez.
    \begin{equation}
        \sum_{i\in I} y_{o,i} \leq 1 \quad \forall o\in O
    \end{equation}
    \item Presupuesto disponible.
    \begin{equation}
        \sum_{i\in I}\left(f_i\,x_i+v_i\sum_{o\in O} y_{o,i}\right)\leq B
    \end{equation}
    \item Compatibilidad técnica observación-instrumento.
    \begin{equation}
        y_{o,i} \leq a_{o,i} \quad \forall i\in I,\ \forall o\in O
    \end{equation}
    \item Naturaleza de las variables.
    \begin{equation}
        y_{o,i},\,x_i \in \{0,1\} \quad \forall i\in I,\ \forall o\in O.
    \end{equation}
\end{itemize}
}
\end{comment}

\section*{Resolución P1}

Se nos solicita minimizar los costos netos de operación y mantenimiento del parque eólico,
que se definen por los (Costos totales de mantenimiento) - (Ingresos obtenidos por generación
de energía). Para esto, es conveniente definir las siguientes variables de decisión:

\begin{table}[htbp]
    \centering
    \small
    \caption{Variables de decisión}
    \label{tab:variables_decision}

    \begin{tabular}{p{0.15\textwidth} p{0.14\textwidth} p{0.55\textwidth}}
        \hline
        \textbf{Variable} & \textbf{Dominio} & \textbf{Descripción} \\
        \hline

        $x_{g,t}$ & $\{0,1\}$ &
        Indica si la turbina $g$ se encuentra en mantenimiento en el período $t$. \\

        $y_{g,t}$ & $\{0,1\}$ &
        Indica si la turbina $g$ se encuentra disponible en el período $t$. \\

        $z_t$ & $\{0,1\}$ &
        Indica si el equipo de mantenimiento se encuentra desplegado en el período $t$. \\

        $p_{g,t}$ & $\mathbb{R}^{+}$ &
        Indica la energía producida por la turbina $g$ en el período $t$. \\

        \hline
    \end{tabular}
\end{table}

Tras esto, podemos plantear la función objetivo del problema, la cual, tal como se mencionó
previamente, consiste en minimizar los costos netos, obteniéndose así que:

\begin{equation}
    FO:
    \min
    \sum_{g \in G}
    \sum_{t \in T}
    C_{g,t} \cdot x_{g,t}
    +
    \sum_{t \in T}
    A_t \cdot z_t
    -
    \sum_{g \in G}
    \sum_{t \in T}
    \pi_t \cdot p_{g,t}
    \tag{1}
\end{equation}

Donde el primer término representa los costos en los que se incurre por la reparación de cada turbina; el segundo término representa el costo asociado a desplegar el equipo de mantenimiento; mientras que el tercer término hace referencia a los ingresos que posee el parque eólico por la venta de energía. Luego, se definen las restricciones del problema, donde se consideran:

\begin{itemize}

    \item \textbf{Reparación única:} Fuerza a cada turbina a recibir mantenimiento en una
    única fecha.

    \begin{equation}
        \sum_{t \in T} x_{g,t} = 1,
        \quad \forall g \in G
    \end{equation}

    \item \textbf{Despliegue del equipo:} Fuerza a que en caso de que una turbina se encuentre
    siendo reparada en el período $t$, se despliegue el equipo de reparación.

    \begin{equation}
        x_{g,t} \leq z_t,
        \quad \forall g \in G,\; t \in T
    \end{equation}

    \item \textbf{Capacidad máxima:} Establece que no se puede calendarizar la reparación de
    más de $M_t$ turbinas en un tiempo $t$.

    \begin{equation}
        \sum_{g \in G} x_{g,t} \leq M_t,
        \quad \forall g \in G,\; t \in T
    \end{equation}

    \item \textbf{Mal clima:} Establece que el equipo solo se puede desplegar cuando las condiciones
    climáticas lo permiten, donde los días $K \subset T$ proveen las condiciones seguras.

    \begin{equation}
        z_t \leq 1 - \mathbf{1}_{\{t \notin \mathcal{K}\}},
        \quad \forall t \in T
    \end{equation}

    \item \textbf{Cumplir demanda:} Representa el requerimimiento mínimo de generación que
    debe tener el parque eólico.

    \begin{equation}
        \sum_{g \in G} p_{g,t} \geq D_t,
        \quad \forall t \in T
    \end{equation}

    \item \textbf{Límites de generación:} Establece que los límites técnicos de las turbinas.

    \begin{equation}
        0 \leq p_{g,t} \leq P_{g,t} \cdot y_{g,t},
        \quad \forall g \in G,\; t \in T
    \end{equation}

    \item \textbf{Disponibilidad:} Representa si las turbinas se encuentran activas o no. Una turbina
    puede estar no activa en el período $t$ si aún no supera el tiempo de falla $\tau_g$ pero se
    encuentra en mantención; o si ya ha superado el tiempo de falla, y aún no ha sido
    reparada.

    \begin{equation}
        y_{g,t} =
        \begin{cases}
            1 - x_{g,t}, & \text{si } t < \tau_g, \\[4pt]
            \displaystyle \sum_{r=1}^{t-1} x_{g,r}, & t \geq \tau_g.
        \end{cases}
    \end{equation}

\end{itemize}